%%% main.tex --- Praxeon: A Praxeological Foundation for Agentic AI
%%%
%%% Portable skeleton: compiles with a stock TeX Live using the base `article`
%%% class plus common packages (no exotic dependencies). To target a specific
%%% venue, swap the class line below:
%%%   ACM:   \documentclass[sigconf,natbib=false]{acmart}
%%%   IEEE:  \documentclass[conference]{IEEEtran}
%%%   LNCS:  \documentclass{llncs}
%%% and adjust the front matter accordingly.
%%%
%%% Build:  latexmk -pdf main.tex   (or `make`)

\documentclass[11pt,a4paper]{article}

\usepackage[utf8]{inputenc}
\usepackage[T1]{fontenc}
\usepackage{lmodern}
\usepackage[margin=1in]{geometry}
\usepackage{microtype}
\usepackage{amsmath,amssymb}
\usepackage{booktabs}
\usepackage{listings}
\usepackage{xcolor}
\usepackage[numbers,sort&compress]{natbib}
\usepackage[hidelinks]{hyperref}

%% A minimal Lisp/Coalton listing style.
\definecolor{lstkw}{RGB}{0,90,160}
\definecolor{lstcm}{RGB}{110,110,110}
\definecolor{lststr}{RGB}{160,60,60}
\lstdefinelanguage{Coalton}{
  morekeywords={coalton-toplevel,define,define-type,define-class,define-struct,
    declare,match,defpackage,defmethod,defgeneric,defclass,defstruct,defun,
    handler-case,restart-case,let,let*,lambda},
  sensitive=true,
  morecomment=[l]{;;},
  morecomment=[l]{;},
  morestring=[b]",
}
\lstset{
  language=Coalton,
  basicstyle=\ttfamily\small,
  keywordstyle=\color{lstkw}\bfseries,
  commentstyle=\color{lstcm}\itshape,
  stringstyle=\color{lststr},
  columns=fullflexible,
  keepspaces=true,
  showstringspaces=false,
  frame=single,
  framesep=4pt,
  breaklines=true,
}

\title{\textbf{Praxeon}: A Praxeological Foundation for\\
       Industrial-Strength Agentic AI in Common Lisp}
\author{Bob\\ \small\texttt{eternal.recursion@proton.me}}
\date{\today}

\begin{document}
\maketitle

\begin{abstract}
The dominance of Python in artificial intelligence is widely attributed to the
language's merits, yet Python's role is largely that of an orchestration layer
over compiled C++/CUDA kernels; the computationally decisive path is already not
Python. We argue that this observation reframes, rather than refutes, the case
for a Lisp in the modern AI stack: the layer a host language can genuinely own is
\emph{agent orchestration and context management}, which is CPU-bound and where
homoiconicity, a condition system, and image-based development are decisive. We
present \textsc{Praxeon}, a Common Lisp framework whose core ontology is drawn
from praxeology---the study of human action---and whose implementation pairs a
statically typed core (in Coalton) with a dynamic shell. We describe a case
study, \emph{Elise}, a psychologically-informed successor to Weizenbaum's ELIZA,
and outline \emph{Kairos}, a bitemporal knowledge substrate for context. We
close with an evaluation plan comparing \textsc{Praxeon} against contemporary
Python and TypeScript agent frameworks on expressiveness and orchestration cost.
\end{abstract}

\section{Introduction}
\label{sec:intro}
% Motivate: PAIP lineage; the frustration with Python's dominance; the claim that
% a Lisp offers a better speed/flexibility trade-off for the orchestration layer.
Since Norvig's \emph{Paradigms of AI Programming}~\citep{norvig1992paip}, Common
Lisp has been a natural medium for expressing artificial-intelligence systems as
programs that manipulate programs. This paper revisits that tradition for the era
of large language models (LLMs), and asks what, precisely, a Lisp can contribute
that the incumbent Python and TypeScript ecosystems cannot.

\section{The Misdiagnosis of Python's Dominance}
\label{sec:misdiagnosis}
% Python as glue over C++/CUDA; the hot path is not Python; therefore no host
% language gets "closer to Rust" on tensors. The real, winnable territory.
Python did not prevail in AI on the strength of the language. Frameworks such as
PyTorch and JAX dispatch the substantive computation to compiled kernels
(cuBLAS, cuDNN, XLA); the interpreter's own throughput is nearly irrelevant
because the workload is kernel-launch- and I/O-bound. A corollary follows: no
choice of host language moves the tensor hot path closer to a systems language,
because that path is already not the host language. The layer that remains
open---and CPU-bound---is orchestration.

\section{A Praxeological Model of Agency}
\label{sec:praxeology}
% ends/means/action/plan/actor; the mapping table; why this beats node/chain.
We model an agent using the vocabulary of praxeology~\citep{mises1949human}: an
\emph{actor} applies \emph{means} to attain \emph{ends} through \emph{action},
ranking ends by preference and economizing a scarce resource. The mapping to
agent architecture is direct (Table~\ref{tab:map}), and, we argue, more
principled than the ``node/chain/runnable'' vocabulary of current frameworks.

\begin{table}[h]
\centering
\caption{Praxeological concepts and their agent-architecture counterparts.}
\label{tab:map}
\begin{tabular}{ll}
\toprule
Praxeology & Agent architecture \\
\midrule
ends & goals \\
means & tools \\
action & applying a means toward an end \\
preference over ends & the objective / valuation \\
action under uncertainty & the LLM's probabilistic reasoning \\
economizing scarce time & the context/token budget \\
imputation of value & credit assignment for tool calls \\
\bottomrule
\end{tabular}
\end{table}

\section{Why a Lisp, and Which}
\label{sec:whichlisp}
% Summarize the dialect analysis: SBCL native perf, Coalton typed core, vs
% Clojure ecosystem, vs Scheme purity. Justify Common Lisp + Coalton.
We compare Common Lisp (SBCL), Clojure, and Scheme (Chez) along native
performance, metaprogramming, static-typing options, and ecosystem, and select
Common Lisp with a Coalton~\citep{coalton} typed core: it offers native
ahead-of-time compilation, the richest macro and condition systems, and---via
Coalton's Hindley--Milner types---a statically checked core suitable for
industrial use.

\section{The Praxeon Architecture}
\label{sec:arch}
% Typed Coalton core (ontology) + dynamic CL shell (effects, conditions, LLM).
\textsc{Praxeon} is organized in two layers. A typed core, written in Coalton,
defines what agents \emph{are}; a dynamic Common Lisp shell owns effects, the
provider protocol, the budgeted context, and the deliberate--act loop. The core
ontology is small:

\begin{lstlisting}
(coalton-toplevel
  (define-type End    (End String))            ; a goal
  (define-type Means  (Means String String))   ; name, description
  (define-type Action (Action Means End String))
  (define-type Plan   (Plan (List Action)))
  (define-type Actor  (Actor String (List Means))))
\end{lstlisting}

Recoverable failure---ubiquitous in long-running agents---is handled with the
condition system, so that recovery policy (retry, substitute, abandon) is
injected at the handler rather than fixed at the point of failure.

\section{Case Study: Elise}
\label{sec:elise}
% The modern-psychology ELIZA; reflective listening via LLM; safety guardrail.
\emph{Elise} is a reflective, psychologically-informed conversational actor---a
successor to Weizenbaum's ELIZA~\citep{weizenbaum1966eliza} that replaces
hard-coded scripts with an LLM prompted to draw on evidence-based practice. Elise
is a demonstration of the actor model, not a clinical instrument, and includes a
transparent safety guardrail that surfaces crisis resources.

\section{Live Development: Reshaping a Running Agent}
\label{sec:live}
% Image-based development in practice: hot reload with state preserved; failures
% surfaced in REPL and browser at once; the structural contrast with Python.
The abstract's claim that image-based development is decisive is best appreciated
in practice. \textsc{Praxeon}'s web surface presents Elise as a server-rendered
(HTMX) application behind a file watcher in the spirit of ClojureScript's
Figwheel~\citep{figwheel}: saving a source file recompiles it into the
\emph{running} image and rebuilds the server around the \emph{same} agent, so the
conversation---the accumulated context---survives while the machinery beneath it
is replaced. Because the request dispatcher invokes its rendering and handler
functions by name, a redefined function is resolved through the function cell on
the next request, and behavior hot-swaps without a restart.

The condition system renders failure legible rather than fatal. A compilation
error during reload is captured and surfaced \emph{simultaneously} as diagnostics
in the developer's REPL and as an overlay on the live page, while the last
successfully compiled server continues to serve; correcting the error clears both
and reloads the view. In use, one edits source and watches the running artifact
change---the conversation intact, errors shown in place---rather than restarting
the process and replaying inputs to recover the prior state.

The contrast with the incumbent stack is structural rather than cosmetic. The
canonical Python loop is to edit, restart the process, lose in-memory state, and
replay inputs to return to the point of interest; auto-reloading facilities
re-execute modules and discard the live object graph, and neither a first-class
means to redefine a function in a running image while preserving state, nor a
condition system to pause and resume at the point of failure, is available. For
agent orchestration in particular---processes that are long-lived, stateful, and
failure-prone---a restart is costly precisely because it discards the scarce
accumulated context. A substrate on which behavior can be reshaped without
discarding that context is therefore not a mere convenience but an \emph{economy},
consonant with this paper's praxeological thesis: the least costly viable means to
a working system. A natural extension makes a runtime condition's restarts
themselves selectable from the browser---the interactive debugger, delivered over
the wire.

\section{Kairos: A Bitemporal Context Substrate}
\label{sec:kairos}
% valid-time vs transaction-time; XTDB-style; Arrow FFI; auditable context.
Context items in \textsc{Praxeon} carry two timestamps---valid time and
transaction time---so that an agent can reason about \emph{what was believed, and
as of when}. We sketch \emph{Kairos}, a standalone bitemporal knowledge graph in
the spirit of XTDB~\citep{xtdb}, intended to back context management via a
foreign-function interface to Apache Arrow for columnar storage.

\section{Related Work}
\label{sec:related}
% Lang* (Python), Mastra (TS), Bosquet/Spell (Clojure); PAIP; Coalton; XTDB.
We situate \textsc{Praxeon} relative to Python agent frameworks, the TypeScript
framework Mastra, and Clojure efforts such as Bosquet, contrasting their
orchestration vocabularies with the praxeological model.

\section{Evaluation Plan}
\label{sec:eval}
% Expressiveness (LOC / concept count for equivalent agents); orchestration
% latency on CPU-bound routing; robustness under injected tool failures.
We propose to evaluate \textsc{Praxeon} on (i) expressiveness, measured by the
concepts and lines required to express equivalent agents; (ii) orchestration
cost, measured as CPU-bound routing and context-assembly latency against
equivalent Python/TypeScript implementations; and (iii) robustness under
injected tool failures, exercising the condition-system recovery paths.

\section{Conclusion}
\label{sec:conclusion}
The winnable claim for a Lisp in modern AI is not raw numeric speed but a
principled, malleable orchestration layer. \textsc{Praxeon} pursues that claim
with a praxeological core, a typed/dynamic split, and a bitemporal context
substrate.

\bibliographystyle{plainnat}
\bibliography{references}

\end{document}
